\begin{figure}[htbp]
\centering
\begin{tikzpicture}
\node[fpblock] (in) at (0,0) {输入 I/O 寄存\quad 1 拍\\RGB888、DE、VS};
\node[fpblock] (lut) at (5,0) {Gamma 寄存 / 同步 LUT\\RGB、DE、VS 共 2 拍};
\node[fpblock] (convert) at (10,0) {RGB888 → RGB565\quad 1 拍\\\texttt{wr\_data / wr\_en / VS}};
\node[fpblock] (stats) at (0,-2.3) {原始亮度与 DE 统计\\$S,N,N_{\rm dark},N_{\rm high}$};
\node[fpblock] (target) at (5,-2.3) {移位 / 加法阈值比较\\连续两帧相同 → 目标曲线};
\node[fpblock] (active) at (10,-2.3) {SOF 锁存使能 / 活动曲线\\向旧目标每帧移动至多一级};
\draw[fparrow] (in)--(lut);
\draw[fparrow] (lut)--(convert);
\draw[fparrow] (in)--(stats);
\draw[fparrow] (stats)--(target);
\draw[fparrow] (target)--(active);
\draw[fparrow,dashed] (active.north)--(10,-1.1)-| (lut.south);
\node[fpnote] at (5,1.0) {关闭 Gamma 选择同延迟旁路；总输入延迟为 4 拍，启动间隔为 1 拍};
\node[fpnote] at (5,-3.35) {统计在下一 SOF 发布，新分类只能影响后续完整帧；首帧使用恒等曲线};
\end{tikzpicture}
\caption{自适应 Gamma 的像素流水线与帧级控制回路}
\label{fig:fpga-gamma}
\end{figure}
